\documentclass{article}

\usepackage[french]{babel}
\usepackage[utf8]{inputenc}
\usepackage[T1]{fontenc}
\usepackage{verbatim}
\usepackage{graphicx}
\usepackage{calc}
\usepackage{color}
\usepackage{float}
\usepackage{hyperref}
\usepackage{geometry}

\geometry{margin=2cm}

\title{Projet système \\ Mise en place d'une bibliothèque de threads}
\author{Benoît Ruelle, Ludovic Hofer}

\begin{document}
\begin{center}
  \includegraphics [width=40mm]{ENSEIRB-MATMECA.jpg} 

\vspace{\stretch{1}}

\textsc{\Huge Simulation de particules sur GPU}\\[0.5cm]
\rule{0.4\textwidth}{1pt}

\vspace{\stretch{1}}

\begin{center}
  
  \begin{flushleft}
    \large
    \emph{Auteurs :}\\
    \begin{itemize}
    \item Benoît Ruelle
    \item Ludovic Hofer
    \end{itemize}
  \end{flushleft}
  
  
  \begin{flushright}
    \large
    \emph{Encadrant :}\\
    Raymond Namyst
  \end{flushright}
\end{center}

\vspace{\stretch{1}}

{\large \url{https://github.com/uOptim/pmg-enseirb/}}

\vspace{\stretch{1}}

{\large Deuxième année, filière informatique}

~

{\large 14 avril 2013 - 6 juin 2013}\\

\end{center}
\thispagestyle{empty}
\pagebreak
\tableofcontents
\newpage

\section{Introduction}

\section{Développement et difficultés rencontrées}

\subsection{Division par zéro dans le cas des collisions}

\paragraph{}
Un des premiers problèmes auxquels nous avons du faire face se situait dans le
noyau \verb!collision!. En effet, lors de l'exécution de la configuration
\verb!bounce.conf!, si la collision était activée dès le début, le
comportement obtenu correspondait à nos attentes, mais si nous laissions les
atomes se traverser une fois avant d'activer la collision, les deux atomes
disparaissaient lors de la collision.

\paragraph{}
En nous penchant plus en détail sur ce problème, nous nous sommes aperçus que le
problème venait du fait qu'à plusieurs endroits de l'équation, on divisait un
élément par $1 + i.x$. Or, dans le cas évoqué ci-dessus, $i.x = -1$, on effectue
donc une division par zéro. Pour pallier ce problème, nous avons simplement
traité ce cas en inversant le vecteur de distance dans ce cas précis.

\subsection{Mise en place de la mesure}
\paragraph{}
Pour mesurer le temps d'exécution, nous avons essayé de mesurer celui-ci en
effectuant la différence entre le temps au début de la fonction idle et à la
fin de celle-ci, ceci afin de mesurer réellement le nombre d'image par seconde.
\paragraph{}
Cette première méthode nous a posé problème étant donné que nous avions un temps
minimal d'environ 16 millisecondes, ce quelque soit les modes et le nombre
d'atomes utilisés. Le temps de 16 millisecondes correspondant étrangement avec
60 images par seconde, nous avons eu un axe pour la recherche du problème, nous
avons rapidement pu déterminer deux sources possibles, la présence d'une
synchronisation verticale et l'utilisation du double buffer. Si nous ne sommes
pas parvenu à enlever la synchronisation verticale, l'utilisation du simple
buffer nous a permis de dépasser les 60 images par secondes et d'avoir ainsi des
mesures plus intéressantes.
\paragraph{}
Étant donné que le rendu graphique n'était que secondaire, nous avons finalement
dérivé sur des mesures de temps avant et après les
\verb!clEnqueueNDRangeKernel!. Afin que la tâche soit réellement mesurée, il
a aussi été nécessaire que nous imposions un \verb!clFinish! avant de prendre
la mesure du temps.

\subsection{Collision v3}
Afin de d'optimiser l'exécution de la fonction collision, nous avons cherché à
rassembler les threads en workgroups afin qu'ils partagent de la mémoire.
\subsubsection{Répartition des workgroups/threads}
tikz pour schéma?
\subsubsection{Calcul de la position en fonction du group id}
\subsubsection{Complexité}
Nombre de threads élevé, mais calcul de la racine carrée... à voir si en
changeant la slice on améliore les résultats
\subsubsection{En faisant varier la slice}
Tests à faire..

\subsection{Lennard Jones v2}

\subsection{Génération de fichiers pour les atomes}
Blabla sur le fichier de génération python : espace grand pour éviter la
diffusion des atomes sur lennard Jones

\subsection{Script de mesure de performances}

\section{Performances mesurées}

\subsection{Cas de la fonction collision}
Graphiques + commentaires
\subsection{Cas du calcul des forces de Lennard-Jones}
Graphiques + commentaires

\section{Conclusion}

\end{document}
